\documentclass[9pt]{book}
\setlength{\oddsidemargin}{-.5 in}
\setlength{\evensidemargin}{-.5 in}
\addtolength{\topmargin}{-1in}
\setlength{\textwidth}{7.5in}
\setlength{\textheight}{9in}
\usepackage{enumitem}
\usepackage{tikz}

\usepackage{amsmath, amssymb}
\usepackage{algorithm}
\usepackage{url, color, epsfig, amsmath, amssymb}
\usepackage{graphicx}
\usepackage{amsmath, amsthm, amssymb}
\usepackage{algpseudocode}

\usepackage{algorithm}
\usepackage{url, color, epsfig, amsmath, amsthm, amssymb}
\usepackage{amsthm}
\usepackage{graphicx}
\usepackage{mathtools}
\DeclarePairedDelimiter{\ceil}{\lceil}{\rceil}
\DeclarePairedDelimiter{\floor}{\lfloor}{\rfloor}

\newcommand{\remove}[1]{}

\newcommand{\Do}{{\small\bf do}\ }
\newcommand{\Proc}[1]{#1\+}
\newcommand{\Returns}{{\small\bf returns}}
\newcommand{\Procbegin}{{\small\bf begin}}
\newcommand{\Then}{{\small\bf then}\ \=\+}
\newcommand{\Elseif}{\<{\small\bf elseif}\ }
\newcommand{\Endif}{\<{\small\bf end\ if\ }\-\-}
\newcommand{\Endproc}[1]{{\small\bf end} #1\-}
\newcommand{\Endfor}{{\small\bf end\ for}\ \-}
\newcommand{\Endloop}{{\bf end\ loop}\ \-}
\newenvironment{program}{
    \begin{minipage}{\textwidth}
        \begin{tabbing}
            \ \ \ \ \=\kill
        }{
        \end{tabbing}
    \end{minipage}
}

\def\blankline{\hbox{}}

\newcommand{\lecture}[5]{
    \pagestyle{headings}
    \thispagestyle{plain}
    \newpage
    \setcounter{chapter}{#1}
    \setcounter{page}{#2}
    \noindent
    \begin{center}
        \framebox{
            \vbox{
                \hbox to 6.28in { {\bf Algorithms
                        \hfill Spring Semester, 2026} }
                \vspace{4mm}
                \hbox to 6.28in { {\Large \hfill Homework 5 (Graded) \hfill} }
                \vspace{2mm}
                \hbox to 6.28in { {\it Professor: #4 \hfill Deadline: 11:59AM Apr 5, 2026.} }
            }
        }
    \end{center}
    \vspace*{4mm}
}

\newcommand{\topic}[2]{\section{#1} \index{#2} \markright{#1}}
\newcommand{\subtopic}[2]{\subsection{#1} \index{#2}}
\newcommand{\subsubtopic}[2]{\subsubsection{#1} \index{#2}}

\renewcommand{\cite}[1]{[#1]}
\renewcommand{\thefootnote}{\arabic{footnote}}

\newcommand{\bi}{\begin{itemize}}
\newcommand{\ei}{\end{itemize}}
\newcommand{\be}{\begin{enumerate}}
\newcommand{\ee}{\end{enumerate}}
\newcommand{\blank}{\vspace{1ex}}

\newtheorem{theorem}{Theorem}[chapter]
\newtheorem{lemma}[theorem]{Lemma}
\newtheorem{claim}[theorem]{Claim}
\newtheorem{digress}[theorem]{Digression}
\newtheorem{corollary}[theorem]{Corollary}

\newenvironment{dfn}{{\vspace*{1ex} \noindent \bf Definition }}{\vspace*{1ex}}
\newcommand{\bigdef}[2]{\index{#1}\begin{dfn} {\rm #2} \end{dfn}}
\newcommand{\smalldef}[1]{\index{#1} {\em #1}}

\begin{document}
    \lecture{0}{1}{\today}{Dr. Mudassir Shabbir}{180 mins.}

    \textbf{Instructions.} This is a graded homework. Submit your solutions as a single PDF. Show all work; correct answers without justification will receive no credit. You may collaborate on high-level ideas but must write up solutions independently.

\bigskip

\begin{enumerate}[label=\arabic*.]

% ---------------------------------------------------------------
% PROBLEM 1 -- LONGEST PALINDROMIC SUBSEQUENCE
% ---------------------------------------------------------------

\item \textbf{[25 points]} \textit{This is an interval DP over substrings.}

Given a string $S=s_1s_2\cdots s_n$, let $\mathrm{LPS}(i,j)$ denote the length of the longest palindromic subsequence in substring $S[i..j]$.

Use the string:
\[
S = \texttt{BANANAS}
\]
(indexed from 1 to 7).

\begin{enumerate}[label=(\alph*)]
    \item \textbf{[5 pts]} Give a precise definition of the state $\mathrm{LPS}(i,j)$ and all base cases.

    \item \textbf{[8 pts]} Write the recurrence for the two cases $s_i=s_j$ and $s_i\ne s_j$.

    \item \textbf{[7 pts]} Compute enough table entries to determine $\mathrm{LPS}(1,7)$. Show your fill order by increasing substring length.

    \item \textbf{[3 pts]} Report the final length of the LPS of \texttt{BANANAS}.

    \item \textbf{[2 pts]} Give asymptotic time and space complexity of the bottom-up DP.
\end{enumerate}

\bigskip

% ---------------------------------------------------------------
% PROBLEM 3 -- PARTITION INTO K SEGMENTS
% ---------------------------------------------------------------

\item \textbf{[25 points]} \textit{You are given an array and must partition it into exactly $k$ contiguous segments to minimize total cost.}

Let $A[1..n]$ be an array of positive integers. For any segment $A[t..i]$, define its segment cost as
\[
\mathrm{cost}(t,i)=\left(\sum_{r=t}^{i} A[r]\right)^2.
\]
You must partition $A[1..n]$ into exactly $k$ nonempty contiguous segments.

For this problem use:
\[
A=[3,1,4,1,5,9],\quad n=6,\quad k=3.
\]

\begin{enumerate}[label=(\alph*)]
    \item \textbf{[5 pts]} Define a DP state $\mathrm{DP}[i][j]$ that represents an optimization over prefixes and segment counts.

    \item \textbf{[5 pts]} Give base cases for $\mathrm{DP}[i][1]$ and any impossible states.

    \item \textbf{[7 pts]} Derive a recurrence for $\mathrm{DP}[i][j]$ by choosing the last segment boundary.

    \item \textbf{[5 pts]} Compute the optimal value for $\mathrm{DP}[6][3]$. Show intermediate values clearly for all states with $j\in\{1,2,3\}$.

    \item \textbf{[3 pts]} State time complexity of the straightforward bottom-up implementation in terms of $n$ and $k$, and explain how prefix sums change the cost of evaluating each transition.
\end{enumerate}

\bigskip

% ---------------------------------------------------------------
% PROBLEM 4 -- INTERLEAVING STRINGS (2D DP)
% ---------------------------------------------------------------

\item \textbf{[25 points]} \textit{Determine if one string can be formed by interleaving two others.}

You are given three strings:
\[
s_1=\texttt{aabcc},\qquad s_2=\texttt{dbbca},\qquad s_3=\texttt{aadbbcbcac}.
\]
% Define $\mathrm{Inter}[i][j]$ = \texttt{true} if $s_3[1..i+j]$ can be formed by interleaving $s_1[1..i]$ and $s_2[1..j]$ in some order (alternating characters, but not necessarily 1-1 alternation).

\begin{enumerate}[label=(\alph*)]
    \item \textbf{[5 pts]} Give a precise state definition and all base cases.

    \item \textbf{[8 pts]} Derive the recurrence. (Hint: consider where the last character of $s_3[i+j]$ came from.)

    \item \textbf{[6 pts]} Fill enough of the DP table to determine $\mathrm{Inter}[5][5]$ (whether $s_3[1..10]$ is an interleaving of $s_1$ and $s_2$).

    \item \textbf{[3 pts]} Report whether $s_3$ can be formed as an interleaving of $s_1$ and $s_2$.

    \item \textbf{[3 pts]} Give asymptotic running time and space complexity.
\end{enumerate}

\bigskip

% ---------------------------------------------------------------
% PROBLEM 5 -- MINIMUM PATH SUM IN GRID (2D DP)
% ---------------------------------------------------------------

\item \textbf{[25 points]} \textit{Find the path from top-left to bottom-right that minimizes total cost.}

You are given a $4 \times 5$ grid with positive integer costs. You may only move right or down.

\[
\begin{array}{ccccc}
1 & 3 & 1 & 5 & 1 \\
2 & 2 & 4 & 1 & 5 \\
5 & 0 & 2 & 3 & 0 \\
6 & 1 & 1 & 2 & 1 \\
\end{array}
\]

% Define $\mathrm{DP}[i][j]$ as the minimum cost to reach cell $(i,j)$ from cell $(1,1)$.

\begin{enumerate}[label=(\alph*)]
    \item \textbf{[5 pts]} Give the state definition and all base cases.

    \item \textbf{[8 pts]} Derive the recurrence relation for $\mathrm{DP}[i][j]$, considering that we can reach $(i,j)$ from either $(i-1,j)$ or $(i,j-1)$.

    \item \textbf{[6 pts]} Fill the complete $\mathrm{DP}$ table for the given grid.

    \item \textbf{[3 pts]} Report the minimum path sum to reach $(4,5)$.

    \item \textbf{[3 pts]} State the time and space complexity, and describe how to reconstruct one optimal path.
\end{enumerate}

\end{enumerate}

\end{document}
